% Shared manuscript body - included by BOTH ral_draft.tex (journal/final,
% IEEEtran) and ral_submission.tex (RA-L initial submission, ieeeconf
% conference format, anonymous). Wrappers define \ifanonymous and the
% paperkeywords environment. Edit content HERE, not in the wrappers.

\begin{abstract}
Soft continuum arms are highly redundant: many actuation patterns realize the
same tip position. Deterministic learned inverse kinematics (IK) regressors
average these incompatible solutions into invalid ones, while a conditional
diffusion model can represent the full solution distribution. We present a
controlled simulation study that asks two questions the existing literature
leaves open: \emph{when} does this multimodality actually help, and does the
advantage survive a change of dynamics model? On a piecewise-constant-curvature
(PCC) pneumatic arm, a conditional diffusion model matches a deterministic
baseline where redundancy is absent and dominates where it is present
(mode recall $0.84\pm0.03$ vs.\ exactly $0$; obstacle-avoidance success
$97.8\%$ vs.\ $4.0\%$; three seeds), and matches a strong optimization-based
solver on accuracy at a tenth of its query cost without needing the
simulator at query time. Two controlled manipulations show the
advantage is specifically a redundancy effect: conditioning on the whole arm
shape instead of the tip collapses sampled diversity by ${\sim}85\times$ and
shrinks the baseline's disadvantage from ${\sim}30\times$ to ${\sim}3\times$,
while increasing the number of segments widens the gap. However, when
actuations from the PCC-trained model are executed on a high-fidelity
Cosserat-rod simulator, accuracy collapses on every seed
($0.67\,$mm $\rightarrow 139.9\,$mm, $0\%$ success): the learned inverse map
is simulator-specific. We characterize what recovers this loss: candidate
selection under a query budget helps but is bounded; a diversity-aware query
order significantly outperforms a confidence order (paired $t{=}3.3$);
guidance-based adaptation yields a modest, statistically suggestive gain;
gain-domain-randomization fails entirely; and fine-tuning on target-domain
data reveals a sharp retention--transfer cliff with no data-mixing sweet
spot. All code, data, and
per-experiment result files are made available.
\end{abstract}

\begin{paperkeywords}
Modeling, Control, and Learning for Soft Robots; Deep Learning Methods;
inverse kinematics; sim-to-sim transfer.
\end{paperkeywords}

\section{Introduction}

\IEEEPARstart{A}{multi-segment} pneumatic continuum arm with three chambers
per segment has more actuation degrees of freedom than task degrees of
freedom: for a four-segment arm, a 12-D actuation vector determines a 3-D tip
position. The inverse map is therefore one-to-many, and its preimages are not
merely a connected null space---distinct arm shapes reach the same tip. A
deterministic learned IK regressor trained by motor
babbling~\cite{giorelli2015neural} fits the conditional \emph{mean} of this
multimodal distribution and can return an actuation that realizes none of the
true solutions. A conditional diffusion
model~\cite{ho2020denoising,song2021denoising} instead learns the full
conditional distribution $p(\mathbf{q}\,|\,\text{target})$ and can sample all
solution modes.

Diffusion models for robot action generation are now
well-established~\cite{chi2023diffusion}, and diffusion-based IK has been
demonstrated for rigid multi-arm systems~\cite{zhang2025ikdiffuser}. What the
literature does not yet provide for \emph{soft continuum} arms is (i) a
controlled account of when the multimodality actually matters---as opposed to
diffusion simply being a strong regressor---and (ii) an honest measurement of
whether the learned inverse map survives the model mismatch that any
simulation-trained soft-arm controller must eventually face. Soft arms make
question (ii) unusually pressing: their dynamics are hard to model, so the
gap between any two models of the \emph{same} physical arm is large, and a
method whose advantage evaporates under that gap is of limited practical
value.

We answer both questions in a fully reproducible sim-to-sim setting: an
idealized, fast PCC model~\cite{webster2010design} as the training simulator,
and a Cosserat-rod simulation (PyElastica~\cite{gazzola2018forward,
zhang2019modeling}) of the same arm as the high-fidelity evaluation domain.
Our contributions:
\begin{enumerate}
  \item \textbf{A validated advantage.} With three-seed statistics, diffusion
  IK dominates learned MLP and mixture-density-network (MDN) baselines,
  recovers $0.84\pm0.03$ of ground-truth IK modes where they structurally
  recover none, and matches a classical random-restart optimization solver
  on accuracy at a tenth of its query cost without simulator access at query
  time---translating into a $97.8\%$ vs.\ $4.0\%$ obstacle-avoidance gap
  over the learned baselines (Sec.~\ref{sec:core}).
  \item \textbf{The mechanism, isolated.} Two controlled manipulations of
  redundancy---conditioning on the whole body shape, and varying the number
  of segments---show the advantage is specifically redundancy resolution,
  not generic model capacity (Sec.~\ref{sec:mechanism}).
  \item \textbf{A quantified transfer collapse---and what recovers it.} The
  learned inverse map collapses under simulator change on every seed. We
  quantify the recovery achievable by inference-time selection and guidance,
  show a diversity-aware query ordering significantly outperforms the obvious
  confidence ordering (paired $t{=}3.3$, $p{\approx}0.002$), show
  gain-domain-randomization achieves nothing, and map a sharp
  retention--transfer cliff for target-domain fine-tuning
  (Sec.~\ref{sec:transfer}).
\end{enumerate}
Negative and bounded results are reported with the same prominence as
positive ones; Sec.~\ref{sec:limitations} states the limitations explicitly.

\section{Related Work}

\textbf{IK for continuum arms.} Non-learned approaches solve continuum IK by
optimization through a model at query time: Jacobian-based schemes for
nonconstant-curvature arms~\cite{giorelli2015neural} and real-time FEM
inversion~\cite{duriez2013control}; the constant-curvature kinematics
underlying our fast simulator are reviewed in~\cite{webster2010design}, and
STIFF-FLOP-style pneumatic arms motivate the three-chamber actuation
model~\cite{cianchetti2014soft}. We include an optimization-based baseline
of this family (random-restart gradient IK through the differentiable
model). Learned approaches fit a deterministic target-to-actuation map on
motor-babbling data~\cite{giorelli2015neural}; mixture density
networks~\cite{bishop1994mixture} are the classical answer to multimodal
regression and serve as the learned multimodal baseline.

\textbf{Diffusion models for robots.} Denoising diffusion
models~\cite{ho2020denoising} with DDIM sampling~\cite{song2021denoising} and
classifier-free guidance~\cite{ho2022classifier} underpin diffusion
policies~\cite{chi2023diffusion}. IKDiffuser~\cite{zhang2025ikdiffuser}
solves IK for \emph{rigid multi-arm} systems with flexible kinematic-tree
\emph{topology}; physically compliant continuum bodies are outside its scope.
Selecting one action from a diffusion policy's samples via learned scorers is
an active area; our transfer-guided sampling instead injects the score's
gradient \emph{during} generation, in the spirit of classifier
guidance~\cite{dhariwal2021diffusion}, and our query-budget framing asks how
many \emph{expensive target-domain evaluations} selection needs, which we
have not seen posed for soft-arm IK.

\textbf{Crossing the reality gap.} Domain
randomization~\cite{tobin2017domain} and target-domain fine-tuning are the
standard levers; catastrophic forgetting under naive fine-tuning is
well-documented~\cite{kirkpatrick2017overcoming}. We evaluate both levers in
a controlled sim-to-sim instance and find the first fails completely while
the second exposes a sharp retention--transfer tension.

\section{Setup}
\label{sec:setup}

\textbf{Arm and training simulator.} The training domain is a four-segment
PCC arm (segment length $60\,$mm, total length $240\,$mm). Each segment has
three $120^\circ$-spaced chambers with normalized pressures
$\mathbf{q}\in[0,1]^{12}$; a linear chamber-to-curvature map (gain $25$)
yields per-segment curvature components $(k_x,k_y)$, and forward kinematics
composes constant-curvature arcs. The model is batched and differentiable.

\textbf{Data and models.} Training data are $5\times10^5$ i.i.d.\ uniform
actuations pushed through forward kinematics. The diffusion model is a
residual MLP denoiser (six FiLM-conditioned blocks, width 512) trained as a
DDPM ($T{=}1000$, cosine schedule) with classifier-free guidance
(condition dropout $0.1$); sampling uses DDIM. Ablations
(Sec.~\ref{sec:core}) fix guidance weight $1.5$ and $20$ DDIM steps.
Baselines: an MLP regressor and an MDN (eight mixture components), matched
in width, plus a classical optimization-based reference---gradient-descent
IK through the differentiable PCC model from 32 random restarts (Adam, 300
iterations), denoted grad-IK. Unless stated otherwise we sample $K{=}32$
candidates per target and report best-of-$K$; success means tip error under
$5\,$mm ($2\%$ of arm length). The obstacle task (E3) runs 200 trials, each placing three random
spherical obstacles (radius $2\,$cm) in the workspace; success requires any
of the $K$ samples to be simultaneously accurate and collision-free.

\textbf{Ground-truth modes.} For evaluation targets we enumerate the true IK
solution set by a $2\times10^6$-sample sweep, keeping actuations whose tip
lands within tolerance and clustering with DBSCAN~\cite{ester1996dbscan}
($\varepsilon{=}12$, $\text{min\_samples}{=}5$; a generated sample matches a
mode within radius $15$). Clustering, mode matching, and diversity are
computed in the 8-D per-segment \emph{curvature} space $(k_x,k_y)$ rather
than raw pressures: the chamber-to-curvature map has a per-segment null
space (common-mode pressure offsets change $\mathbf{q}$ but not the shape),
which would otherwise contaminate distance-based metrics. Evaluation targets
average $2.05$ ground-truth modes (range 0--6); targets where the sweep
finds no mode are excluded from recall averages.

\textbf{High-fidelity domain.} The transfer domain is a Cosserat-rod
simulation of the same arm in PyElastica~\cite{gazzola2018forward,
zhang2019modeling}: one rod, 48 elements, per-segment constant internal
couples derived from the same chamber pressures, damped to quasi-static
equilibrium (${\sim}4\,$s of compute per evaluation). The two simulators
agree exactly in the zero-curvature case---uniform pressures produce the
identical straight configuration (tip at $[0,0,0.24]\,$m) in both---so the
discrepancy lies in the bending response, not in trivial misalignment. A
sweep over the couple gain further confirmed it is directional/structural
rather than a scale mismatch: no gain setting aligns the two workspaces.

\section{When Multimodality Helps}
\label{sec:core}

\begin{table}[t]
\caption{Core comparison on the PCC arm (mean $\pm$ std, 3 seeds; $K{=}32$).
Mode recall and diversity are computed over successful samples only.
grad-IK is a single optimization run per target (no training, hence no
seeds) and requires the differentiable simulator at query time.}
\label{tab:core}
\centering
\scriptsize
\setlength{\tabcolsep}{2.2pt}
\begin{tabular}{lcccc}
\toprule
 & Diffusion & grad-IK & MLP & MDN \\
\midrule
Tip err.\ best-of-$K$ (mm) & $0.44\pm0.03$ & $\mathbf{0.23}$ & $13.4\pm0.0$ & $17.3\pm0.2$ \\
Success @ 5 mm & $\mathbf{1.00}$ & $0.994$ & $0.060$ & $0.027$ \\
Mode recall & $\mathbf{0.84\pm0.03}$ & $0.73$ & $0.00$ & $0.00$ \\
Diversity (curv.\ space) & $\mathbf{28.2\pm0.2}$ & $9.5$ & $0.0$ & $0.0$ \\
Obstacle success & $\mathbf{0.978\pm0.002}$ & $0.960$ & $0.040$ & $0.018$ \\
Query time (ms/target) & $38.6\pm4.9$ & $395$ & $\mathbf{0.013}$ & $0.022$ \\
\bottomrule
\end{tabular}
\end{table}

Table~\ref{tab:core} summarizes the four core experiments. \textbf{E1
(accuracy):} best-of-$K$ diffusion reaches $0.44\,$mm mean tip error with
$100\%$ success; the deterministic MLP, structurally limited to the
conditional mean, plateaus at $13.4\,$mm, and the MDN does not improve on
it. The classical grad-IK reference reaches $0.23\,$mm ($99.4\%$)---when the
true model is available to optimize through at query time, classical
optimization is fully competitive on accuracy.
\textbf{E2 (multimodality):} against enumerated ground-truth solution sets,
diffusion recovers $0.84\pm0.03$ of modes; both baselines recover none---the
MLP because it is deterministic, and the MDN not through classic component
collapse but because no component is accurate enough: mode recall and
diversity are computed over \emph{successful} samples only, and the MDN
almost never produces one. Per-component analysis rules out a training
pathology---NLL decreases monotonically at every epoch, no variance reaches
its clamp floor, and $7.8$ of $8$ components remain effective---yet even the
\emph{best} component mean per target carries $7.5\,$mm error (median
component $9.2\,$mm) against the $5\,$mm tolerance, and sampling with the
learned $\sigma$ smears this further. Nor is it a tuning artifact: sweeping
the component count over $\{4,8,16,32\}$ improves best-of-$K$ error only
from $18.7$ to $14.8\,$mm---never surpassing the plain MLP---and mode recall
remains exactly zero at every setting. grad-IK's random restarts, by
contrast, \emph{do} recover modes (recall $0.73$)---multimodality does not
require learning---though with a third of diffusion's solution diversity
($9.5$ vs.\ $28.2$). \textbf{E3 (downstream value):} when success requires \emph{any} of
$K$ samples to be both accurate and collision-free among random obstacles,
diffusion succeeds on $97.8\%$ of trials vs.\ $4.0\%$ (MLP)---redundancy is
only exploitable if the model can represent it. \textbf{E4 (cost):}
$38.6\,$ms per target for a batch of 32 samples on CPU, versus $395\,$ms
for grad-IK---and, more fundamentally, grad-IK requires the differentiable
simulator in the loop at every query, whereas the learned methods answer
from samples of the plant alone. That amortization property is exactly what
Sec.~\ref{sec:transfer} puts under stress. With three seeds,
formal hypothesis testing adds little; the separations are instead
unambiguous by inspection: on every metric in Table~\ref{tab:core} the
\emph{worst} diffusion seed outperforms the \emph{best} baseline seed---for
tip error, $0.49$ vs.\ $13.33\,$mm, a gap exceeding $400\times$ the
across-seed standard deviation.

Three ablations support the protocol: guidance weight has an inverted-U
optimum at $1.5$ (both accuracy and diversity collapse beyond ${\sim}3$);
accuracy and recall saturate by ${\sim}20$ DDIM steps (a $3\times$ speedup
over 50 steps at no measurable cost); and performance saturates by
${\sim}0.5$--$1\times10^5$ training samples (Fig.~\ref{fig:scaling}).

\begin{figure}[t]
\centering
\includegraphics[width=\linewidth]{fig_scaling.png}
\caption{Dataset-size ablation (single seed). (a) Best-of-$K$ tip error and
(b) mode recall vs.\ training-set size; both saturate by
${\sim}0.5$--$1\times10^5$ samples.}
\label{fig:scaling}
\end{figure}

\section{The Advantage Is Redundancy, Specifically}
\label{sec:mechanism}

\begin{figure*}[t]
\centering
\includegraphics[width=0.95\textwidth]{fig_redundancy.png}
\caption{The diffusion advantage tracks actuation redundancy. (a) Replacing
the 3-D tip target with a 24-D whole-shape target removes most of the
redundancy: sampled diversity collapses ${\sim}85\times$ ($28.16\pm0.22$ to
$0.346\pm0.022$; mean $\pm$ std, 3 seeds). (b) Increasing the number of
segments (single seed per morphology) widens the diffusion--MLP accuracy gap;
at six segments the obstacle-task gap reaches $100\%$ vs.\ $1\%$.}
\label{fig:redundancy}
\end{figure*}

If the advantage in Table~\ref{tab:core} is redundancy resolution, it should
shrink when redundancy is removed and grow when redundancy is added. Both
predictions hold (Fig.~\ref{fig:redundancy}).

\textbf{Removing redundancy: shape conditioning.} Conditioning the same
models on the full 24-D body shape (eight backbone points) instead of the
3-D tip pins down almost the entire configuration. Diffusion's sampled
diversity collapses from $28.16\pm0.22$ to $0.346\pm0.022$
(${\sim}85\times$), and the MLP's accuracy disadvantage shrinks from
${\sim}30\times$ ($13.4$ vs.\ $0.44\,$mm) to ${\sim}3\times$ ($0.49$ vs.\
$0.16\,$mm), with both models at $100\%$ success (3 seeds). When there is
one answer, a deterministic regressor is nearly as good; diffusion's edge
is specifically the multimodal case.

\textbf{Adding redundancy: morphology.} Training fresh models on two- and
six-segment variants of the same arm family (tolerance scaled to $2\%$ of
arm length), the MLP's error grows from $1.7\,$mm (2 seg, $75\%$ success) to
$33.9\,$mm (6 seg, $1.2\%$ success) while diffusion degrades gracefully
($0.11$ to $0.92\,$mm, ${\geq}99.8\%$ success); the obstacle-task gap at six
segments is $100\%$ vs.\ $1\%$---the largest in the study. More redundancy
means worse mode averaging and more collision-free alternatives that only a
multimodal model can offer.

\section{Does It Survive Model Mismatch?}
\label{sec:transfer}

\subsection{The transfer collapse}
Executing the PCC-trained diffusion model's actuations on the Cosserat-rod
simulator collapses accuracy from $0.67\pm0.03\,$mm ($100\%$ success) to
$139.9\pm1.8\,$mm with $0\%$ success on every seed ($n{=}20$ held-out
targets, $K{=}8$); the MLP collapses similarly ($13.8 \rightarrow
164.4\,$mm). The learned inverse map is simulator-specific. The rest of this
section quantifies what does---and does not---recover the loss;
Table~\ref{tab:synthesis} summarizes.

\begin{table}[t]
\caption{Recovery mechanisms at a glance (Elastica-domain performance;
best configuration of each mechanism).}
\label{tab:synthesis}
\centering
\footnotesize
\setlength{\tabcolsep}{2.2pt}
\begin{tabular}{lccl}
\toprule
Mechanism & err (mm) & succ. & Cost / caveat \\
\midrule
none (PCC-trained) & $139.9\pm1.8$ & $0\%$ & --- \\
best-of-32 selection & $130.4$ & $0\%$ & 32 target queries/target \\
\; + diversity order & $-6.7$ at $m{=}8$ & $0\%$ & same queries; $p{\approx}0.002$ \\
guided sampling & $129.0$ & $0\%$ & 400 labels; suggestive \\
gain randomization & $140.6$ & $0\%$ & degrades source; no gain \\
fine-tune, $0\%$ mix & $\mathbf{67.1\pm1.0}$ & $\mathbf{10\%}$ & forgets source ($133\,$mm) \\
fine-tune, $\geq 5\%$ mix & $\geq 133.8$ & $\leq 10\%$ & retains source, forfeits gain \\
\bottomrule
\end{tabular}
\end{table}

\subsection{Selection under a query budget}
\begin{figure}[t]
\centering
\includegraphics[width=\linewidth]{fig_query_budget.png}
\caption{Candidate selection under an Elastica query budget ($n{=}45$
targets, $K{=}32$ candidates, mean $\pm$ SEM). Curves shown are independent
runs of the two orderings; the statistical comparison in the text uses a
\emph{paired} design (one candidate set per target, both orderings reindex
the same cached per-candidate errors), which is significant at $m{=}8$ and
$m{=}16$ ($t{=}3.34$ and $3.21$, $p{\approx}0.002$).}
\label{fig:budget}
\end{figure}

A multimodal model offers something a deterministic one cannot: $K$ diverse
candidates to test against the target domain. Spending a budget of $m$
Elastica evaluations and keeping the best candidate reduces mean error from
$161.5$ ($m{=}1$) to $130.4\,$mm ($m{=}32$), improving 44 of 45 targets
(Fig.~\ref{fig:budget})---real and broad-based, but to be blunt about its
practical value: success stays at $0\%$, the gap to tolerance remains
${\sim}26\times$, and under this budget selection never actually solves a
single target. \emph{How} the
budget is spent matters: a greedy farthest-point ordering in curvature space
(seeded from the same most-confident candidate, so $m{=}1$ is identical)
beats confidence ordering by $6.7\,$mm at $m{=}8$ (paired $t{=}3.34$,
$p{\approx}0.002$; 24/12 targets) and $3.3\,$mm at $m{=}16$ ($t{=}3.21$).
Notably, neither candidate diversity itself nor local PCC Jacobian
sensitivity predicts \emph{which} targets benefit: at $n{=}45$---sized to
detect $r{\geq}0.41$ at $80\%$ power, i.e., the effect size an $n{=}20$
pilot had suggested---the observed correlations are $r{=}0.20$ ($p{=}0.19$)
and $r{=}0.02$ ($p{=}0.88$). This rules out the pilot's effect but not
smaller ones; the pilot's non-replication is a caution we report rather
than hide.

\subsection{Guiding generation toward the target domain}
Beyond selecting among candidates, we steer generation itself: a small
regressor trained on 400 Elastica-labeled candidates (25 targets) predicts
transfer error from (candidate, target), and its gradient nudges the
denoised estimate at every DDIM step, as in classifier
guidance~\cite{dhariwal2021diffusion}. On 15 \emph{held-out} targets this
improves best-of-16 error from $134.6$ to $129.0\,$mm (paired
$5.6\pm2.5\,$mm, $t{=}2.27$, $p{\approx}0.04$; 11/15 targets). This is
generalization to unseen targets, but it is the weakest effect in the
paper: all $p$-values here are uncorrected, and across the family of
${\sim}$six hypothesis tests we report, only the ordering result
($p{\approx}0.002$) survives a Bonferroni correction---we therefore read
this one as \emph{suggestive} rather than confirmed.
Combining it with diversity-ordering adds nothing: guidance
\emph{concentrates} the candidate distribution while diversity-ordering
\emph{spreads} selection, and a $2\times2$ factorial shows the mechanisms in
tension rather than compounding.

\subsection{Training-time levers}
\textbf{Domain randomization fails.} Randomizing the curvature gain
per-sample ($\mathcal{U}(15,35)$) during data generation costs PCC accuracy
($0.44 \rightarrow 2.18\,$mm) and yields \emph{zero} transfer improvement
($140.6 \rightarrow 140.6\,$mm)---consistent with the mismatch being
directional rather than magnitude-based: randomizing how much the arm bends
cannot fix a discrepancy in which direction it bends.

\textbf{Target-domain data works---at a price.}
\begin{table}[t]
\caption{Fine-tuning on $10^4$ Elastica-simulated samples at six PCC-data
mixing ratios (mean $\pm$ std, 3 seeds; held-out targets).}
\label{tab:cliff}
\centering
\footnotesize
\setlength{\tabcolsep}{3.5pt}
\begin{tabular}{lcccc}
\toprule
PCC mix & PCC err (mm) & succ. & Elastica err (mm) & succ. \\
\midrule
none (baseline) & $0.67\pm0.03$ & $100\%$ & $139.9\pm1.8$ & $0\%$ \\
$0\%$  & $133.4\pm1.2$ & $0\%$ & $\mathbf{67.1\pm1.0}$ & $\mathbf{10\%}$ \\
$5\%$  & $30.7\pm0.6$ & $6\%$ & $133.8\pm3.1$ & $8\%$ \\
$10\%$ & $21.7\pm0.5$ & $17\%$ & $143.2\pm4.9$ & $8\%$ \\
$15\%$ & $16.7\pm0.2$ & $26\%$ & $146.0\pm3.4$ & $8\%$ \\
$20\%$ & $11.6\pm0.3$ & $51\%$ & $148.0\pm1.6$ & $10\%$ \\
$50\%$ & $3.0\pm0.4$ & $96\%$ & $152.3\pm0.4$ & $3\%$ \\
\bottomrule
\end{tabular}
\end{table}

\begin{figure}[t]
\centering
\includegraphics[width=\linewidth]{fig_finetune_cliff.png}
\caption{The retention--transfer cliff (mean $\pm$ std, 3 seeds; six mixing
ratios). Fine-tuning the PCC-pretrained model on $10^4$ Elastica samples:
retention (PCC error) improves smoothly and monotonically as PCC data is
mixed back in, but the transfer gain collapses almost entirely by the first
grid point ($5\%$ PCC data). There is no data-mixing sweet spot.}
\label{fig:cliff}
\end{figure}

Fine-tuning the pretrained model \emph{exclusively} on $10^4$
Elastica-simulated samples yields the only nonzero transfer success of any
method---$67.1\pm1.0\,$mm and $10\%$ success ($2/20$ targets on each seed;
recall the $5\%$ metric granularity)---at the cost of catastrophic
forgetting~\cite{kirkpatrick2017overcoming} of the source domain ($0.67
\rightarrow 133.4\,$mm). Mixing PCC data back into fine-tuning batches
restores retention smoothly and monotonically across all six ratios
($133.4 \rightarrow 30.7 \rightarrow 21.7 \rightarrow 16.7 \rightarrow 11.6
\rightarrow 3.0\,$mm), but the transfer gain does not degrade in step: the
fine sweep \emph{locates} the collapse below the first grid point---at just
$5\%$ PCC data, transfer error has already rebounded from $67.1$ to
$133.8\pm3.1\,$mm, forfeiting ${\sim}92\%$ of the improvement over the
no-fine-tune baseline, and it drifts only slowly thereafter
(Table~\ref{tab:cliff}, Fig.~\ref{fig:cliff}). The values at intermediate
ratios are stable across seeds (std $\leq 4.9\,$mm), so this is not an
unlucky training run. Retention and transfer sit on a \emph{cliff} whose
edge lies below $5\%$ mixing, not on a smooth tradeoff curve---under
per-batch data mixing there is no sweet spot, suggesting parameter-level
regularization (\eg\ elastic weight
consolidation~\cite{kirkpatrick2017overcoming}) as the necessary next
mechanism. (Binary success has $5\%$ granularity at
$n_\text{targets}{=}20$; mean error carries the statistical weight
throughout.)

\section{Limitations}
\label{sec:limitations}
(1) \textbf{Simulation-only.} Both domains are simulators; this is a
controlled sim-to-sim study, and no claim is made about physical hardware.
(2) \textbf{Simplified target-domain actuation.} The Cosserat-rod arm applies
idealized per-segment couples and quasi-static settling; it is a harder test
domain, not a validated model of a specific physical arm. (3)
\textbf{Statistical coverage is uneven.} Core results, the shape-conditioning
study, the baseline transfer collapse, and the fine-tuning grid carry
three-seed statistics; the query-budget, ordering, and guidance experiments
use one trained model (with within-experiment pairing across 15--45 targets),
the hyperparameter, dataset-size, and morphology studies are likewise
single-seed, and the grad-IK reference is a single optimization run per
target.
(4) \textbf{Mode-enumeration hyperparameters} were calibrated for the
four-segment arm; morphology results therefore report accuracy and
obstacle-task metrics only. (5) \textbf{Success-rate granularity} in transfer
experiments is $5\%$ ($n_\text{targets}{=}20$).

\section{Conclusion}
For the PCC-representable pneumatic continuum arms studied here, multimodal
diffusion IK delivers a real, mechanistically attributable advantage---it is
redundancy resolution, growing with segment count and vanishing when the
target specifies the whole shape---but the learned inverse map itself does
not survive a change of dynamics model.
Inference-time selection and guidance recover a bounded fraction of the loss;
synthetic gain randomization recovers none; and target-domain fine-tuning
buys genuine transfer only at a sharp retention cliff that data mixing cannot
smooth. For simulation-trained soft-arm IK, the practical path forward is not
better sampling from a source-domain model but mechanisms that let a small
amount of target-domain data reshape the model without erasing it.

\section*{Reproducibility}
\ifanonymous
All code, configuration, per-experiment result files, and the
$10^4$-sample Elastica dataset are available in an anonymized repository:
\url{https://anonymous.4open.science/r/REPLACE-BEFORE-SUBMISSION}.
\else
All code, configuration, per-experiment result files, and the
$10^4$-sample Elastica dataset are available at
\url{https://github.com/adhvayjagadeesh/soft-arm-diffusion-ik}.
\fi
